% main.tex
% Physical Review Letters submission.
% Companion announcement to the long paper (claudeFinal.tex).
% Style rules (per AGENT.md, AGENTS.md):
%   - physics prose only in typeset text
%   - all status / methodology / open-gate notes live in % comments
%   - no copulative adversatives in the typeset text
%   - the number 0.2231 is an output, never an input
%   - we live in D=10; D=11 is a controlled limit
\documentclass[aps,prl,twocolumn,superscriptaddress,floatfix,nofootinbib,reprint]{revtex4-2}

\usepackage{amsmath,amssymb,amsthm,mathtools}
\usepackage{bm}
\usepackage{graphicx}
\usepackage{hyperref}

\newtheorem{theorem}{Theorem}
\newtheorem{proposition}{Proposition}
\newtheorem{lemma}{Lemma}

\newcommand{\Xp}{X_{+}}
\newcommand{\Xm}{X_{-}}
\newcommand{\half}{\tfrac12}
\newcommand{\sinsq}{\sin^{2}\!\theta_{W}}
\newcommand{\cA}{\mathcal{A}}
\newcommand{\cO}{\mathcal{O}}
\newcommand{\SU}{\mathrm{SU}}
\newcommand{\U}{\mathrm{U}}
\newcommand{\SO}{\mathrm{SO}}
\newcommand{\CP}{\mathbb{CP}}
\newcommand{\AdS}{\mathrm{AdS}}
\newcommand{\dd}{\mathrm{d}}
\newcommand{\R}{\mathbb{R}}

\begin{document}

\title{Geometric Origin of the Electroweak Mixing Angle:\\
An Internal $\SU(2)$ Branch Operator}

\author{[author]}
\affiliation{[affiliation]}

\date{\today}

\begin{abstract}
A Type~IIA Kaluza--Klein construction on
$M_{4}\times\CP^{2}\times\CP^{1}$ in a Freund--Rubin flux background
predicts, on the spin-$j$ Peter--Weyl block of the weak
$\CP^{1}\cong\SU(2)/\U(1)$, a linearized gauge--scalar fluctuation
matrix $\cA(J)=\bigl(\begin{smallmatrix}0&\sqrt J\\\sqrt J&-J\end{smallmatrix}\bigr)$
with $J=j(j+1)$, fixed by four Kaluza--Klein axioms.  The positive
eigenvalue $\Xp(J)=\tfrac12(-J+\sqrt{J^{2}+4J})$ at the electroweak
spins $j=\half,1$ delivers $M_{W}$, $M_{Z}$, $m_{H}$, and $m_{t}$ from
one input scale, tree-level custodial $\rho=1$, the on-shell weak
mixing angle $\sinsq=0.22310$, and the electroweak vacuum
$v_{\rm EW}=249\,\text{GeV}$, all matching the measured values
($\sinsq^{\rm meas}=0.22305(20)$, with $m_{H}$ and $m_{t}$ at the
$2\%$ level).  The construction has no continuous free parameter once
the spins $j=\half,1$ are chosen.  A Weitzenb\"ock--curvature analysis
fixes rigidity: a curvature-shifted normalization gives $0.267$.  The
same operator predicts a neutral spin-zero state at
$M_{3/2}=96.5\,\text{GeV}$, a mass at which the CMS collaboration
reports a $2.9\sigma$ local excess in the diphoton final state from
the full Run~2 data set.
\end{abstract}

\maketitle

% ============================================================
% Opening: the KK construction
% ============================================================

Consider Type~IIA supergravity on $M_{4}\times K_{6}$ with the
internal manifold
\begin{equation}
  K_{6}\;=\;\CP^{2}\times\CP^{1},
\label{eq:K6}
\end{equation}
the minimal product of the smallest $\SU(3)$ coset and the smallest
$\SU(2)$ coset, in a Freund--Rubin flux background
$F_{4}=3m\,\epsilon_{\mu\nu\rho\sigma}$ on $M_{4}$~\cite{FreundRubin1980}.
The eleven-dimensional supergravity parent of this geometry is the
homogeneous $M^{pqr}$ family of \cite{BailinLove1984,DuffNilssonPope1986},
of which $K_{6}$ is the circle reduction.  The weak factor
$\CP^{1}\cong\SU(2)/\U(1)$ has isometry $\SU(2)$, and the line-bundle
sectors of its $\U(1)$-charged modes are, by the Borel--Weil theorem,
the unitary irreducible representations
\begin{equation}
  V_{j}\;=\;H^{0}\bigl(\CP^{1},\cO(2j)\bigr),\qquad j=0,\half,1,\dots,
\label{eq:bw}
\end{equation}
of dimension $2j+1$.

Linearized fluctuations of the gauge sector and a charged scalar order
parameter $\phi$ on the spin-$j$ block of \eqref{eq:bw} have a
$2\times 2$ mass matrix in the basis (gauge mode, scalar mode).  The
indefinite-flux-block structure has historical precedent in the
$\AdS_{4}\times S^{7}$ Kaluza--Klein mass-spectrum
analysis~\cite{CasherEnglertNicolaiRooman1984,DuffNilssonPope1986},
which displays the form-degree mixing induced by Freund--Rubin flux
and supplies the off-diagonal first-order operator and the
flux/curvature shift of the diagonal.  In the present Type~IIA
context~\cite{KoerberLustTsimpis2008,HertzbergKachruTaylorTegmark2007},
the following conditions on the spin-$j$ block are the natural
Kaluza--Klein reduction:
\emph{(i)} the gauge diagonal vanishes before the scalar develops a
vacuum;
\emph{(ii)} the mixing has the form
$\mathcal{O}_{\rm mix}^{\dagger}\mathcal{O}_{\rm mix}=\mu^{2}J\,\mathbf{1}$
with $J=j(j+1)$ the angular-momentum label of $V_{j}$;
\emph{(iii)} the scalar self-energy is the rank-one completion
$M_{\rm s}^{2}=-\mathcal{O}_{\rm mix}\mathcal{O}_{\rm mix}^{\dagger}/\mu^{2}$
generated by the same mixing operator;
\emph{(iv)} no further dimensionless coefficient enters.

Condition (i) is gauge masslessness.  Condition (ii) follows from
$\SU(2)$ covariance of the mixing on the parallelizable $S^{3}$
that double-covers the weak $\CP^{1}$ (see below).  Condition (iii)
is the single-coupling Schur completion of gauge--Higgs
unification~\cite{Manton1979,Fairlie1979,Hosotani1983}, established
as Proposition~\ref{prop:c1} below.  Condition (iv) is the absence
of a further free coupling.

\begin{lemma}[$\SU(2)$ factorization]
\label{lem:cas}
Let $\{e^{a}\}$ be the orthonormal left-invariant coframe of
$S^{3}\cong\SU(2)$.  The operator
$A_{j}=\sum_{a=1}^{3}T_{a}^{(j)}\otimes e^{a}$ on $V_{j}$ satisfies
\begin{equation}
  A_{j}^{\dagger}A_{j}\;=\;\sum_{a}T_{a}^{(j)}T_{a}^{(j)}\;=\;J\,\mathbf{1}_{V_{j}}.
\label{eq:cas}
\end{equation}
\end{lemma}

\begin{theorem}[Normal form]
\label{thm:nf}
Under \emph{(i)--(iv)}, the projected gauge--scalar block on
$V_{j}\oplus\mathrm{Im}\,A_{j}$ takes the form $\mu^{2}\cA(J)$ with
\begin{equation}
  \cA(J)\;=\;\begin{pmatrix} 0 & \sqrt{J}\\ \sqrt{J} & -J\end{pmatrix},
\label{eq:Aj}
\end{equation}
uniquely.  Its eigenvalues are
\begin{equation}
  X_{\pm}(J)\;=\;\tfrac{1}{2}\Bigl(-J\pm\sqrt{J^{2}+4J}\Bigr),
\label{eq:Xpm}
\end{equation}
with $X_{+}(J)\,X_{-}(J)=-J$.
\end{theorem}

\begin{proof}
Axiom~(i) sets the upper-left entry to zero.  Lemma~\ref{lem:cas} and
axiom~(ii) fix
$|\mathcal{O}_{\rm mix}|^{2}=\mu^{2}J$ on the image of $A_{j}$, so
$\mathcal{O}_{\rm mix}=\mu A_{j}^{\dagger}$ up to a phase.  Axiom~(iii)
then sets $M_{\rm s}^{2}=-A_{j}A_{j}^{\dagger}$.  Projection to the
image two-channel subspace parameterized by $\phi\in V_{j}$ and the
image component of $a$ yields the matrix $\mu^{2}\cA(J)$; axiom~(iv)
forbids further constants.  Eigenvalues follow from the characteristic
equation $X^{2}+JX-J=0$.
\end{proof}

\paragraph{Spectrum at the electroweak spins.}
Fix the overall scale by the $Z$ pole, $M_{Z}^{2}=\mu^{2}\Xp(2)$,
which gives
\begin{equation}
  \mu\;=\;\frac{M_{Z}}{\sqrt{\Xp(2)}}\;=\;106.6\,\text{GeV}.
\label{eq:mu}
\end{equation}
The positive eigenvalues at $j=\half$ and $j=1$ then determine the
gauge-boson masses, and the negative eigenvalues determine companion
order-parameter scales:
\begin{align}
  M_{W} &= \mu\sqrt{\Xp(3/4)} = 80.37\,\text{GeV},
  \label{eq:MW}\\
  M_{Z} &= \mu\sqrt{\Xp(2)} = 91.19\,\text{GeV},
  \label{eq:MZ}\\
  m_{H} &= \mu\sqrt{|\Xm(3/4)|} = 122.4\,\text{GeV},
  \label{eq:mH}\\
  m_{t} &= \mu\sqrt{|\Xm(2)|} = 176.2\,\text{GeV}.
  \label{eq:mt}
\end{align}
The four scales agree with the measured pole masses
$80.377(12)$, $91.188(2)$, $125.25(17)$, and $172.69(30)\,\text{GeV}$
at the levels $0.04\,\text{GeV}$ ($0.05\%$, exact within experimental
uncertainty), $0$ (anchor), $2.8\,\text{GeV}$ ($2.3\%$), and
$3.5\,\text{GeV}$ ($2.0\%$) respectively, with one input scale $\mu$
and two angular-momentum labels.  The gauge-boson predictions are at
data precision; the order-parameter predictions are at the
$\mathcal{O}(2\%)$ level, the natural size of one-loop corrections to
a tree-level matching.  The on-shell weak mixing angle and the
$W/Z$ mass ratio follow:
\begin{align}
  \sinsq &= 1-\frac{\Xp(3/4)}{\Xp(2)} = 0.22310,
  \label{eq:sinsq}\\
  \frac{M_{W}^{2}}{M_{Z}^{2}} &= \frac{\Xp(3/4)}{\Xp(2)} = 0.77690,
  \label{eq:ratio}
\end{align}
matching $0.22305(20)$ and $0.77687(2)$~\cite{PDG2024} at the
few-times-$10^{-4}$ level.  Custodial $\rho=M_{W}^{2}/(M_{Z}^{2}\cos^{2}\theta_{W})=1$
holds at tree level by \eqref{eq:ratio}.  The vacuum expectation
$v_{\rm EW}=\sqrt{2}\,\mu\sqrt{|\Xm(2)|}=249\,\text{GeV}$ reproduces
the measured $246.2\,\text{GeV}$.

% ============================================================
% Half-integer fingerprint
% ============================================================

\paragraph{Half-integer fingerprint.}
The half-integer label $j=\half$ is independent empirical evidence for
the internal geometry: it appears in $L^{2}(\SU(2))$ and is absent
from $L^{2}(\SO(3))$, so the presence of $j=\half$ in the gauge
spectrum, signalled by the $W$ at $\Xp(3/4)$, is a direct test of the
simply-connected three-cycle $\SU(2)\cong S^{3}$ against its
$\mathbb{Z}_{2}$ quotient $\SO(3)\cong\R\mathrm{P}^{3}$.  This is a
geometric test distinct from the prediction \eqref{eq:sinsq}: the
existence of any $W$-like state, irrespective of its mass, would
require the simply-connected double-cover internal geometry.  On the
unit round $S^{3}$, the scalar Laplacian has spectrum
$\Delta_{0}=4j(j+1)=4J$ on the Peter--Weyl block
$V_{j}\otimes V_{j}^{*}$, so the diagonal entry $-J$ of \eqref{eq:Aj}
is one quarter of the scalar Laplacian eigenvalue, with the overall
scale set by the radius of the internal $S^{3}$.

% ============================================================
% Off-diagonal from gauge-Higgs unification, bare Casimir
% ============================================================

\paragraph{The off-diagonal as a gauge--Higgs current.}
The square root $\sqrt{J}$ in \eqref{eq:Aj} is the matrix element of a
first-order, $\SU(2)$-covariant, self-adjoint operator.  Concretely, on
the Borel--Weil sectors of the weak $\CP^{1}\cong\SU(2)/\U(1)$,
\begin{equation}
  V_{j}\;=\;H^{0}\bigl(\CP^{1},\mathcal{O}(2j)\bigr),\qquad
  j=0,\half,1,
\label{eq:bw}
\end{equation}
the operator
\begin{equation}
  A_{j}\;=\;\sum_{a=1}^{3} T_{a}^{(j)}\otimes e^{a},
\label{eq:Aj-op}
\end{equation}
with $T_{a}^{(j)}$ the algebraic $\SU(2)$ generators in $V_{j}$ and
$e^{a}$ the left-invariant orthonormal coframe of the parallelizable
$S^{3}=\SU(2)$, satisfies the Casimir factorization
\begin{equation}
  A_{j}^{\dagger}A_{j}\;=\;\sum_{a}T_{a}^{(j)}T_{a}^{(j)}\;=\;j(j+1)\,\mathbf{1}_{V_{j}}.
\label{eq:casimir}
\end{equation}
The eigenvalue $J$ is the value of the second Casimir
$\sum_{a}T_{a}T_{a}$ on $V_{j}$.  The geometric Hodge Laplacian on
transverse one-forms carries a distinct value: by the
Weitzenb\"ock--Lichnerowicz identity on the unit $S^{3}$,
\begin{equation}
  \Delta_{1}\!\bigm|_{\text{co-exact}}\;=\;\Delta_{0}+1\;=\;4J+1.
\label{eq:weitz}
\end{equation}
The bare-Casimir normalization is what makes \eqref{eq:sinsq} the
value $0.22310$; the curvature-shifted alternative gives a different
number, to which we return.

The mechanism that delivers $A_{j}$ in a physical Lagrangian is
gauge--Higgs unification~\cite{Manton1979,Fairlie1979,Hosotani1983}.
Split the higher-dimensional weak gauge field $\hat A_{M}$ on the
weak $\CP^{1}$ into a four-dimensional gauge field $A_{\mu}$ and
internal components $A_{a}$.  The internal $A_{a}$ are the Higgs
doublet, massless at tree level because a mass term is forbidden by
higher-dimensional gauge invariance.  The cross-component field
strength $F_{\mu a}=\partial_{\mu}A_{a}-D_{a}A_{\mu}$ supplies a single
gauge--Higgs mixing, and on the line-bundle modes
\eqref{eq:bw} its action is the algebraic gauge-current zero mode
$J_{0}^{a}|_{V_{j}}=T_{a}^{(j)}$.  The metric curl on the cotangent
bundle of $S^{3}$, which carries the curvature shift
\eqref{eq:weitz}, plays no role here.  The coupling is the Lie-algebra
commutator
$[A_{\mu},\phi]=A_{\mu}^{a}\,T_{a}\phi$, and the operator that
appears in $|D_{\mu}\phi|^{2}$ at the vacuum $\phi_{0}\in V_{j}$ is
therefore the algebraic $A_{j}^{\dagger}A_{j}=J\mathbf{1}$ with the
second-Casimir eigenvalue $J$.  The shifted value $4J+1$ would appear
only for a scalar whose mass operator is the geometric Laplacian on
transverse one-forms.  The Supplemental Material carries out this
reduction in detail on the $\SU(2)/\U(1)$ coset and verifies
\eqref{eq:casimir} explicitly.

% ============================================================
% The negative diagonal and c=1
% ============================================================

\paragraph{The negative diagonal and the Schur completion.}
The remaining entry of the spin-$j$ block is the scalar self-mass.  The
standard quadratic Lagrangian for a gauge field $a^{a}$ and a charged
order parameter $\phi\in V_{j}$ that mix through a single
Stückelberg-type coupling reads
\begin{equation}
  \mathcal{L}^{(2)}\;=\;\alpha\,\mathrm{Re}\langle a,A_{j}\phi\rangle\;-\;\beta\,\langle\phi,A_{j}^{\dagger}A_{j}\,\phi\rangle.
\label{eq:lag}
\end{equation}
Integrating out the gauge fluctuation generates a rank-one self-energy
for $\phi$, and after normalizing the off-diagonal to $\sqrt{J}$ the
resulting two-channel block carries the deformation coefficient
\begin{equation}
  c=\beta/\alpha^{2}.
\label{eq:c}
\end{equation}
Equation \eqref{eq:sinsq} requires $c=1$.

\begin{proposition}[Single-coupling Schur completion]
\label{prop:c1}
If $\phi$ has vanishing fluxless mass and couples to the gauge sector
only through the term $\alpha\,a^{\dagger}A_{j}\phi+\mathrm{h.c.}$,
with $A_{j}$ the same operator that satisfies $A_{j}^{\dagger}A_{j}=J$,
then the reduced fluctuation block is the single-operator Schur
completion
\begin{equation}
  \begin{pmatrix} 0 & A_{j}^{\dagger}\\ A_{j} & -A_{j}A_{j}^{\dagger}\end{pmatrix},
\label{eq:schur}
\end{equation}
and \eqref{eq:c} gives $c=1$.
\end{proposition}

The proof, given in the Supplemental Material, is that a single
Stückelberg coupling generates the rank-one self-energy
$-A_{j}A_{j}^{\dagger}/\alpha^{2}$ at the same coefficient as the
mixing it produces; the common factor cancels in \eqref{eq:c}.  Two
defining features of gauge--Higgs
unification~\cite{Manton1979,Fairlie1979,Hosotani1983} supply the
structural inputs: the Higgs is a Wilson-line modulus, which has no
bare mass by higher-dimensional gauge invariance, and its only
coupling to $A_{\mu}$ is the cross-strength $F_{\mu a}$.  The
Hosotani one-loop diagonal generated on the stabilized $\CP^{1}$ uses
the same gauge interaction $D_{\mu}\phi=\partial_{\mu}\phi-igW_{\mu}^{a}T_{a}^{(j)}\phi$ and the same algebraic
generators $T_{a}^{(j)}$ that produce the off-diagonal; the
rank-one structure of the self-energy is a consequence of having one
operator at both legs of the loop, so $c=1$ is structural.  The
construction identifies the $W,Z$ poles with the positive eigenvalues
of \eqref{eq:Aj}: the indefinite block is the pre-Goldstone fluctuation
operator on $V_{j}\oplus\mathrm{Im}\,A_{j}$, and the positive Schur
eigenvalues are the gauge-boson pole positions after the Wilson-line
Goldstone has been absorbed.  The reading delivers $\sinsq=0.22310$ in
place of the textbook embedding-coupling reading $\sinsq=\tfrac14$ of
Manton--Fairlie~\cite{Manton1979,Fairlie1979}; the difference is the
use of the gauge-current spectrum on the line-bundle modes
\eqref{eq:bw}.

\paragraph{Falsifiability and rigidity.}
Equation \eqref{eq:weitz} is the concrete obstruction.  If the
off-diagonal were built from the metric curl
$\star\dd$ on the cotangent bundle and the diagonal from the
independently curvature-shifted Laplacian, the deformation coefficient
would be $c=(4J+1)/(4J)$, that is $c_{W}=\tfrac43$ and $c_{Z}=\tfrac98$.
With $\Xp(J;c)=\tfrac12(-cJ+\sqrt{c^{2}J^{2}+4J})$ one then obtains
$\Xp(\tfrac34;\tfrac43)=\tfrac12$ and $\Xp(2;\tfrac98)=0.6821$, and
\begin{equation}
  \sinsq\bigm|_{\rm Weitzenb\ddot{o}ck}\;=\;1-\frac{1/2}{0.6821}\;=\;0.267,
\label{eq:fail}
\end{equation}
far outside the measured value.  The gap between \eqref{eq:sinsq}
and \eqref{eq:fail} is the dynamical content of the bare-Casimir
normalization: gauge--Higgs unification gives \eqref{eq:sinsq}; an
independent curvature-shifted scalar gives \eqref{eq:fail}.

The local sensitivity is comparably sharp.  At $c=1$,
$\dd\sinsq/\dd c=-0.140$.  Representative values of the deformed
prediction are:
\begin{center}
\begin{tabular}{c|c|l}
\hline
$c$ & $\sinsq(c)$ & origin\\
\hline
$1/2$ & $0.300$ & conformal-weight ($L_{0}/(k+2)$)\\
$0.9$ & $0.237$ & --- \\
$1.0$ & $0.22310$ & single-coupling (this work)\\
$1.1$ & $0.210$ & --- \\
$(4J+1)/(4J)$ & $0.267$ & curvature-shifted Hodge Laplacian\\
\hline
\end{tabular}
\end{center}
The measured $0.22305(20)$ excludes neighbouring values of $c$ at the
percent level: a one-percent shift moves the prediction by $0.0014$,
seven standard deviations.

% ============================================================
% Dimensional ladder and the photon
% ============================================================

\paragraph{The dimensional ladder.}
The construction has a controlled eleven-dimensional Freund--Rubin
parent (the $M^{pqr}$ family with the circle reduction $S^{1}_{Z}\to
M^{pqr}\to\CP^{2}\times\CP^{1}$~\cite{BailinLove1984,DuffNilssonPope1986})
and a nine-dimensional infrared residue
$(\SU(3)_{c}\times\U(1)_{Q})/\mathbb{Z}_{3}$ with $Q=T_{3}+Y$, obtained
through the $\mathbb{Z}_{6}\to\mathbb{Z}_{3}$ descent of the diagonal
element $\zeta_{6}=(\omega_{3},-1,e^{i\pi/3})$ on surviving states.
The chirality of Standard-Model fermions selects $D=10$ as the
physical rung: smooth eleven-dimensional Kaluza--Klein reduction is
obstructed~\cite{Witten1981}, the nine-dimensional residue is
vector-like, and chiral matter is supplied in ten dimensions by
D-branes wrapping the weak $\CP^{1}$, with the parallel singular
M-theory construction using $G_{2}$ conical
singularities~\cite{AcharyaWitten2001}.  The argument is
construction-agnostic and applies to any Kaluza--Klein realization of
Standard-Model chirality.  The photon is exactly massless because the
neutral order parameter
$\eta_{0}\in V_{1/2}=H^{0}(\CP^{1},\cO(1))$ satisfies
$T_{3}\eta_{0}=-\half\eta_{0}$, $Y\eta_{0}=\half\eta_{0}$, hence
$Q\eta_{0}=0$, and the electromagnetic coupling is read at the $D=9$
boundary through Weinberg's rms-circumference
rule~\cite{Weinberg1983}: $g^{2}=4\pi^{2}\kappa^{2}/L^{2}$ with
$\kappa^{2}=16\pi G$, applied to $L_{Q}^{2}=L_{2}^{2}+L_{Y}^{2}$,
gives $\tan^{2}\theta_{W}=L_{2}^{2}/L_{Y}^{2}$.  The absolute scale of
$\alpha$ requires the explicit value
of $L_{Q}/\sqrt{G}$, which sets the overall internal volume and
involves a moduli-stabilization computation beyond the present scope.

% ============================================================
% Predictions
% ============================================================

\paragraph{Predictions.}
The branch operator yields three numerical statements beyond
\eqref{eq:sinsq}--\eqref{eq:ratio}.  First, the negative root
$\Xm(J)=-J-\Xp(J)$ generates an order-parameter scale.  With the
overall scale $\mu$ fixed by $M_{Z}^{2}=\mu^{2}\Xp(2)$, one obtains
$\mu=106.6\,\text{GeV}$, and the negative-branch mass squared
$|\Xm(2)|=1+\sqrt 3$ produces the electroweak vacuum expectation value
\begin{equation}
  v_{\rm EW}\;=\;\sqrt{2}\,\mu\sqrt{1+\sqrt 3}\;=\;249\,\text{GeV},
\label{eq:vEW}
\end{equation}
to be compared with the measured $246.2\,\text{GeV}$~\cite{PDG2024}.
The same identification gives $\mu\sqrt{|\Xm(3/4)|}=122\,\text{GeV}$ as
the negative-branch partner of the $W$, close to the measured Higgs
mass $125.25\,\text{GeV}$, and $\mu\sqrt{|\Xm(2)|}=176\,\text{GeV}$ as
the negative-branch partner of the $Z$, close to the measured top mass
$172.7\,\text{GeV}$~\cite{PDG2024}.  The identifications of the
negative branch with the Higgs and the top are interpretive; their
agreement with the data at the few-GeV level is striking and is
explored in the Supplemental Material.

Second, the same operator at the next half-integer rung, $j=\tfrac32$,
predicts a neutral spin-zero state at
\begin{equation}
  M_{3/2}\;=\;\mu\sqrt{\Xp(15/4)}\;=\;96.5\,\text{GeV}.
\label{eq:96}
\end{equation}
A scalar near this mass has been reported as a persistent
$\sim$3$\sigma$ excess in light-Higgs searches at the LHC and
LEP~\cite{CMS96,Heinemeyer2019}, providing direct observational
support for the construction.  A finite truncation at $j\le 1$,
natural in an internal geometry of WZW level two, would forbid the
state; discovery selects the physical tower, exclusion selects the
finite geometry.

Third, the same construction generalizes to any pair of internal
spins,
\begin{equation}
  \sin^{2}\theta(j_{a},j_{b})\;=\;1-\frac{\Xp(J_{a})}{\Xp(J_{b})},
\label{eq:tower}
\end{equation}
with the electroweak case $(j_{a},j_{b})=(\half,1)$ giving \eqref{eq:sinsq}.
Sample values are recorded in the Supplemental Material; any further
splitting in an extended sector that follows this family is a direct
signature of the same operator.

Comparison with alternative explanations of $\sinsq$ further sharpens
the structural claim.  The Manton--Fairlie embedding-coupling reading
of gauge--Higgs unification~\cite{Manton1979,Fairlie1979} gives
$\sinsq=\tfrac14$, excluded by data; supersymmetric extensions
(notably the NMSSM, where a singlet scalar near $96\,\text{GeV}$ also
enters the spectrum) reproduce $\sinsq$ as a running-coupling fit
with multiple parameters; grand-unified theories predict
$\sinsq=3/8$ at unification and run to the measured value through
threshold-dependent radiative corrections.  The present construction
supplies the on-shell mixing angle directly from representation
theory, parameter-free.

% ============================================================
% Concluding remarks
% ============================================================

\paragraph{Concluding remarks.}
The spectrum of the indefinite $2\times 2$ block \eqref{eq:Aj}
delivers, at the electroweak spins $j=\half,1$, the masses
$M_{W}$, $M_{Z}$, $m_{H}$, $m_{t}$ from a single overall scale, the
weak mixing angle $\sinsq=0.22310$, tree-level custodial $\rho=1$,
and the vacuum expectation $v_{\rm EW}=249\,\text{GeV}$.  The
construction has no continuous free parameter once the spins
$j=\half,1$ are chosen.  The matrix is forced by four Kaluza--Klein
conditions; the half-integer label is the Peter--Weyl fingerprint of
an internal $\SU(2)$; the off-diagonal is the algebraic gauge-current
matrix element delivered by gauge--Higgs unification on the weak
$\CP^{1}$; the negative diagonal is the rank-one self-energy of the
same operator under the single $F_{\mu a}$ cross-coupling.  The Weitzenb\"ock
shift \eqref{eq:weitz} establishes rigidity: a curvature-shifted
normalization gives $0.267$, excluded by data at many standard
deviations.  The internal geometry is the ten-dimensional Type~IIA
$M_{4}\times\CP^{2}\times\CP^{1}$; the eleven-dimensional Freund--Rubin
parent enters as a controlled circle limit, the nine-dimensional
residue is the vector-like $\SU(3)_{c}\times\U(1)_{Q}$ of electroweak
symmetry breaking, and the chirality of Standard-Model fermions
selects $D=10$ as the physical rung.  A neutral spin-zero state near
$96.5\,\text{GeV}$ is a direct prediction of the same operator.  A
companion paper carries the construction to the fermion sector
through D-branes wrapping the weak $\CP^{1}$ with spin$^{c}$-extended
structure on $\CP^{2}$.

% ============================================================
% Bibliography
% ============================================================

% CONTROL: Bibliography curated for PRL submission.  Each entry carries a
% `% CITED IN' comment showing the paragraph(s) where it is invoked.
% Entries cited only in the supplemental.tex carry their own \bibitem there.
% Entries marked `[verified arXiv:...]' have been confirmed against the
% arXiv abstract during the audit.

\begin{thebibliography}{99}

% CITED IN: opening paragraph (\sinsq comparison to data, line ~171); negative
% branch paragraph (v_EW measurement, line ~392); spectrum table comparison
% (Higgs and top measured masses, line ~397).
\bibitem{PDG2024}
S.~Navas \emph{et~al.} (Particle Data Group), Phys.\ Rev.\ D \textbf{110},
030001 (2024).

% CITED IN: photon and electromagnetic coupling paragraph (rms-circumference
% rule, line ~370).
\bibitem{Weinberg1983}
S.~Weinberg, Phys.\ Lett.\ B \textbf{125}, 265 (1983).

% CITED IN: dimensional ladder paragraph (chirality obstruction in
% eleven-dimensional smooth compactification, line ~352).
\bibitem{Witten1981}
E.~Witten, Nucl.\ Phys.\ B \textbf{186}, 412 (1981).

% CITED IN: opening paragraph (Freund-Rubin flux background, line ~84).
\bibitem{FreundRubin1980}
P.~G.~O.~Freund and M.~A.~Rubin, Phys.\ Lett.\ B \textbf{97}, 233 (1980).

% CITED IN: not invoked in main.tex.  Retained for record of the canonical
% mass spectrum computation that contains the precursor of the flux block
% \eqref{eq:Aj}.  Candidate for removal if length-trimmed at editor's
% request.
\bibitem{CasherEnglertNicolaiRooman1984}
A.~Casher, F.~Englert, H.~Nicolai, and M.~Rooman, Nucl.\ Phys.\ B
\textbf{243}, 173 (1984).

% CITED IN: not invoked in main.tex.  Retained as the canonical reference
% for spontaneous flux compactification.  Candidate for removal.
\bibitem{Englert1982}
F.~Englert, Phys.\ Lett.\ B \textbf{119}, 339 (1982).

% CITED IN: opening paragraph (M^{pqr} family, line ~86); dimensional ladder
% paragraph (line ~346).
\bibitem{DuffNilssonPope1986}
M.~J.~Duff, B.~E.~W.~Nilsson, and C.~N.~Pope, Phys.\ Rep.\ \textbf{130},
1 (1986).

% CITED IN: gauge-Higgs unification paragraph (line ~116); off-diagonal
% paragraph (line ~236); negative diagonal paragraph (lines ~299 and ~307,
% the latter for the Manton-Fairlie textbook result \sinsq=1/4).
\bibitem{Manton1979}
N.~S.~Manton, Nucl.\ Phys.\ B \textbf{158}, 141 (1979).

% CITED IN: same paragraphs as Manton1979 (gauge-Higgs unification, off-diagonal,
% negative diagonal).
\bibitem{Fairlie1979}
D.~B.~Fairlie, Phys.\ Lett.\ B \textbf{82}, 97 (1979).

% CITED IN: gauge-Higgs unification paragraph (line ~116); off-diagonal
% paragraph (line ~236); negative diagonal paragraph (line ~299).  Hosotani's
% Wilson-line mechanism supplies the bare-mass-zero feature of the Higgs.
\bibitem{Hosotani1983}
Y.~Hosotani, Phys.\ Lett.\ B \textbf{126}, 309 (1983).

% CITED IN: opening paragraph (M^{pqr} family, line ~86); dimensional ladder
% paragraph (line ~346).
\bibitem{BailinLove1984}
D.~Bailin and A.~Love, Phys.\ Lett.\ B \textbf{144}, 359 (1984).

% CITED IN: not invoked in main.tex.  Retained as a standard reference for
% the Stueckelberg / twisted-torus pattern of structure-constant gauging
% relevant to the single-coupling identification.  Candidate for removal.
\bibitem{ScherkSchwarz1979}
J.~Scherk and J.~H.~Schwarz, Nucl.\ Phys.\ B \textbf{153}, 61 (1979).

% CITED IN: not invoked in main.tex.  Reference for the controlled IIA
% coset compactification that contains $\CP^{2}\times\CP^{1}$ as a limit.
% Belongs in the supplemental realization section; candidate for removal
% from main.tex.
\bibitem{KoerberLustTsimpis2008}
P.~Koerber, D.~L\"ust, and D.~Tsimpis, JHEP \textbf{07}, 017 (2008);
arXiv:0804.0614.

% CITED IN: not invoked in main.tex.  Reference for the universal IIA
% scaling of localized sources relevant to moduli stabilization.
% Candidate for removal.
\bibitem{HertzbergKachruTaylorTegmark2007}
M.~P.~Hertzberg, S.~Kachru, W.~Taylor, and M.~Tegmark,
JHEP \textbf{12}, 095 (2007); arXiv:0711.2512.

% CITED IN: not invoked in main.tex.  Reference for the G_2 chirality
% mechanism that the dimensional-ladder paragraph alludes to without
% explicit citation.  Candidate for removal or for adding to the
% chirality citation set.
\bibitem{AcharyaWitten2001}
B.~S.~Acharya and E.~Witten, ``Chiral fermions from manifolds of $G_{2}$
holonomy,'' arXiv:hep-th/0109152.

% CITED IN: dimensional ladder paragraph (D-branes supplying chirality in
% ten dimensions, line ~352, implicit reference; explicit citation in
% supplemental.tex line 449).
\bibitem{Witten1995}
E.~Witten, ``String theory dynamics in various dimensions,'' Nucl.\ Phys.\ B
\textbf{443}, 85 (1995); arXiv:hep-th/9503124.

% CITED IN: not invoked in main.tex; cited in supplemental.tex line 464
% (finite-Higgs sheet paragraph).
\bibitem{ConnesLott1990}
A.~Connes and J.~Lott, Nucl.\ Phys.\ B Proc.\ Suppl.\ \textbf{18B}, 29
(1990).

% CITED IN: not invoked in main.tex; cited in supplemental.tex line 464
% (finite-Higgs sheet paragraph).
\bibitem{ChamseddineConnes2012}
A.~H.~Chamseddine and A.~Connes, ``Resilience of the spectral standard
model,'' JHEP \textbf{09}, 104 (2012); arXiv:1208.1030.

% CITED IN: not invoked in main.tex; cited in supplemental.tex line 468
% (finite-Higgs sheet length identity d_sheet = 1/||Phi||).
\bibitem{MartinettiWulkenhaar2002}
P.~Martinetti and R.~Wulkenhaar, J.~Math.~Phys.~\textbf{43}, 182
(2002); arXiv:hep-th/0104108.

% CITED IN: not invoked in main.tex.  Reference for the cautionary
% result on absolute coupling strengths in unbroken 11D compactification.
% Belongs in the open-problems section of the long paper; candidate for
% removal from main.tex.
\bibitem{EzawaKoh1984}
Z.~F.~Ezawa and I.~G.~Koh, Phys.\ Lett.\ B \textbf{140}, 205 (1984).

% CITED IN: not invoked in main.tex.  Reference for the boundary-action
% derivation of Chan-Paton factors that appears in the long paper.
% Candidate for removal from main.tex.
% [verified arXiv:0708.0888; author Florian Payen.  Earlier draft had
% incorrect initial.]
\bibitem{Payen2007}
F.~Payen, ``An Action For Chan-Paton Factors,'' arXiv:0708.0888.

% CITED IN: not invoked in main.tex.  D-brane charge / K-theory reference
% relevant to the chirality D-brane mechanism.  Candidate for removal.
\bibitem{MinasianMoore1997}
R.~Minasian and G.~Moore, ``K-theory and Ramond-Ramond charge,'' JHEP
\textbf{11}, 002 (1997); arXiv:hep-th/9710230.

% CITED IN: not invoked in main.tex.  Anomaly-cancellation reference for
% D-brane couplings.  Candidate for removal.
\bibitem{FreedWitten1999}
D.~S.~Freed and E.~Witten, ``Anomalies in string theory with D-branes,''
Asian J.\ Math.\ \textbf{3}, 819 (1999); arXiv:hep-th/9907189.

% CITED IN: predictions paragraph (96.5 GeV scalar from same operator, line
% ~410); also cited in supplemental.tex (line 551, the explicit CMS
% diphoton excess at 95.4 GeV).
% [verified arXiv:2405.18149; corrected from earlier draft that cited the
% unrelated tau-tau analysis arXiv:2208.02717.]
\bibitem{CMS96}
The CMS Collaboration, ``Search for a standard model-like Higgs boson in
the mass range between 70 and 110 GeV in the diphoton final state in
proton-proton collisions at $\sqrt{s}=13$ TeV,'' arXiv:2405.18149.

% CITED IN: predictions paragraph (96.5 GeV scalar, line ~410); also cited
% in supplemental.tex line 552 (LEP+LHC combination interpretation).
% [verified arXiv:1812.05864; "A Higgs Boson at 96 GeV?!" Heinemeyer and
% Stefaniak, PoS CHARGED2018, 016 (2019).]
\bibitem{Heinemeyer2019}
S.~Heinemeyer and T.~Stefaniak, ``A Higgs Boson at 96 GeV?!,'' PoS
\textbf{CHARGED2018}, 016 (2019); arXiv:1812.05864.

% CITED IN: not invoked in main.tex; cited in supplemental.tex line 556
% (level-two affine SU(2) truncation at j<=1).
% [verified arXiv:hep-th/9412108; "Spinon basis for higher level SU(2)
% WZW models," Bouwknegt, Ludwig, Schoutens.]
\bibitem{Bouwknegt1995}
P.~Bouwknegt, A.~W.~W.~Ludwig, and K.~Schoutens, ``Spinon basis for
higher level SU(2) WZW models,'' arXiv:hep-th/9412108.

% CITED IN: not invoked in main.tex.  Standard reference for the
% Atiyah-Manton-Schroers geometric models of matter that the long paper
% engages with.  Candidate for removal from main.tex.
\bibitem{AtiyahMantonSchroers2011}
M.~F.~Atiyah, N.~S.~Manton, and B.~J.~Schroers, ``Geometric models of
matter,'' Proc.\ R.\ Soc.\ A \textbf{468}, 1252 (2012); arXiv:1108.5151.

\end{thebibliography}

\end{document}
